\chapter{Concept Selection}
\label{ch:concept_selection}
\section{Session Overview}

Session \ref{ch:concept_generation} generated several rough concepts. This session narrows those ideas to the most promising options using structured comparison tools rather than intuition alone.

By the end of this session, you should be able to:
\begin{enumerate}
    \item derive useful selection criteria from customer needs, specifications, and enterprise considerations;
    \item build a concept-screening matrix;
    \item interpret screening results and justify a shortlist; and
    \item decide whether more detailed concept scoring is necessary.
\end{enumerate}

\section{Key Terms}

\begin{itemize}
    \item \textbf{Selection criterion}: an attribute used to compare concepts, such as ease of maintenance or cost.
    \item \textbf{Concept-screening matrix}: a quick comparison tool that rates concepts relative to a datum concept using symbols such as +, 0, and --.
    \item \textbf{Concept scoring}: a more detailed evaluation method that uses weights and numerical ratings.
    \item \textbf{Datum concept}: the reference concept used for comparison in a screening matrix.
\end{itemize}

\section{Activity 1: Generate Selection Criteria}

% Start by reviewing the outputs from Sessions 7, 8, and 10. Useful criteria usually come from:
% \begin{itemize}
%     \item customer needs, such as easy monitoring or low maintenance;
%     \item product specifications, such as low leakage or stable flow; and
%     \item enterprise concerns, such as cost, manufacturability, or ease of assembly.
% \end{itemize}
By this time, it is expected that you have already familiarized yourself with the chosen product and its specifications.
 \begin{enumerate}
        \item Take note of any specific customer needs which are relevant to the product.
        \begin{enumerate}
            \item Refer to the customer needs (Session 5 activity) that have been identified for the chosen product.
        \end{enumerate}
        \item Take note of any specific enterprise needs which are relevant to the product.
        \begin{enumerate}
            \item Refer to the enterprise needs that have been identified for the chosen product, if available.
        \end{enumerate}

    \end{enumerate}

\subsection{Draft the selection criteria}
        \begin{enumerate}
            \item Choose at least one dimension that will serve as selection criteria for the concepts.
            \item The selection criteria should effectively distinguish among the concepts.
            \item Expressed the criteria at a high level of abstraction.
            \item Do not include unimportant criteria in the screening matrix to avoid difficulty in interpreting results.
        \end{enumerate}

For hydroponic concepts, typical criteria might include:
\begin{itemize}
    \item nutrient delivery performance;
    \item ease of cleaning;
    \item reliability;
    \item product cost;
    \item ease of installation; and
    \item physics-related suitability, such as sensor accuracy or structural stability.
\end{itemize}

\section{Build the Concept-Screening Matrix}

Next, we will set up a concept-screening matrix which is critical when evaluating and comparing different concepts based on specific dimensions.


Lets start:
\begin{enumerate}
    \item Gather the identified concepts:
    \begin{enumerate}
        \item Referring to the concepts that were identified during the Session 7 activity.
    \end{enumerate}
    \item Determine the dimensions for the concept-screening matrix:
    \begin{enumerate}
        \item Refer to criteria selected on activity 1.
    \end{enumerate}
    \item Set up the concept-screening matrix:
    \begin{enumerate}
        \item Take a blank A4 paper and orient it horizontally (landscape).
        \item Draw a table with
        \begin{itemize}
            \item columns for the identified concepts and
            \item rows for the selected dimensions.
        \end{itemize}
        \item Label the
        % \begin{itemize}
        %     \item first column as "Concepts" and write the names or labels of the identified concepts in the subsequent rows.
        %     \item the first row with the selected dimensions.
        % \end{itemize}
    \end{enumerate}
\end{enumerate}


Choose one concept as the datum, then compare the other concepts against it using:
\begin{itemize}
    \item \textbf{+} if the concept is better than the datum on that criterion;
    \item \textbf{0} if it is about the same; and
    \item \textbf{--} if it is worse.
\end{itemize}

Table~\ref{tab:session11-screening-matrix} gives a simple example.

\begin{table}[htbp]
\centering
\caption{Example concept-screening matrix for hydroponic concepts}
\label{tab:session11-screening-matrix}
\renewcommand{\arraystretch}{1.2}
\begin{tabularx}{\textwidth}{|>{\raggedright\arraybackslash}p{0.28\textwidth}|c|c|c|}
\hline
\textbf{Criterion} & \textbf{Concept A (datum)} & \textbf{Concept B} & \textbf{Concept C} \\
\hline
Easy monitoring & 0 & + & 0 \\
\hline
Stable nutrient flow & 0 & 0 & + \\
\hline
Ease of cleaning & 0 & + & -- \\
\hline
Estimated cost & 0 & -- & + \\
\hline
Ease of installation & 0 & 0 & + \\
\hline
\end{tabularx}
\end{table}

After filling in the matrix, count the number of pluses and minuses for each concept. This helps the team see which concepts are consistently stronger.

\section{In-Class Activity 1: Consolidate the Team Matrix}

In class, compare the draft matrices prepared by each team member and produce one shared screening matrix.

\begin{enumerate}
    \item compare the criteria proposed by different members;
    \item remove criteria that do not meaningfully distinguish concepts;
    \item agree on the datum concept;
    \item discuss disagreements in the ratings; and
    \item finalise the screening matrix with a short written rationale.
\end{enumerate}

\section{In-Class Activity 2: Decide Whether to Use Concept Scoring}

Concept scoring is useful when the screening results are too close or when the team needs a more defensible final comparison.

If concept scoring is required, use the following steps:
\begin{enumerate}
    \item keep the key criteria from the screening stage;
    \item assign a weight to each criterion based on its importance;
    \item rate each concept on a numerical scale, such as 1--5; and
    \item multiply ratings by weights to obtain weighted totals.
\end{enumerate}

\begin{table}[htbp]
\centering
\caption{Optional concept-scoring template}
\label{tab:session11-scoring-template}
\renewcommand{\arraystretch}{1.2}
\begin{tabularx}{\textwidth}{|>{\raggedright\arraybackslash}p{0.24\textwidth}|c|c|c|c|}
\hline
\textbf{Criterion} & \textbf{Weight} & \textbf{Concept A} & \textbf{Concept B} & \textbf{Concept C} \\
\hline
Easy monitoring &  &  &  &  \\
\hline
Stable nutrient flow &  &  &  &  \\
\hline
Ease of cleaning &  &  &  &  \\
\hline
Estimated cost &  &  &  &  \\
\hline
Weighted total &  &  &  &  \\
\hline
\end{tabularx}
\end{table}

If your group decides not to use scoring, record the reason clearly. For example, the screening matrix may already show a strong separation between concepts.

\section{Summary and Next Steps}

At the end of this session, your group should have a documented shortlist and a justified reason for the final selection method used. The selected concept or concepts will move into concept testing and refinement in Session 12.
